% Problemas del capítulo 14. Generado por tools/problemas/gen.py; no editar a mano.
\section{Problemas}\label{sec:problemas-14}

Graduados: \(\star\) directo, \(\star\star\) para pensar, \(\star\star\star\) desafío. Las respuestas están en el apéndice~\ref{ap:respuestas}; conviene intentarlos antes de mirarlas.

\begin{problema}[\({\star}\)]\label{pb:14-1}
Para el elipsoide GRS80 (\(a=6\,378\,137\)~m, \(1/f=298{,}257222101\)), calculá el semieje menor \(b\) y la excentricidad al cuadrado \(e^2=2f-f^2\).
\end{problema}

\begin{problema}[\({\star}{\star}\)]\label{pb:14-2}
Calculá los radios de curvatura \(M\) (en el meridiano) y \(N\) (en el primer vertical) en la latitud de Mendoza, \(\varphi=-32{,}89^\circ\). ¿En qué dirección es más curva la Tierra ahí?
\end{problema}

\begin{problema}[\({\star}{\star}\)]\label{pb:14-3}
¿Cuánto mide un grado de latitud en el ecuador y a \(60^\circ\) de latitud? (Un grado de meridiano mide \(M\,\pi/180\).)
\end{problema}

\begin{problema}[\({\star}{\star}\)]\label{pb:14-4}
En Ushuaia (\(\varphi=-54{,}8^\circ\)), calculá el radio del paralelo (\(N\cos\varphi\)) y cuánto mide un grado de longitud. ¿Qué fracción es de lo que mide en el ecuador?
\end{problema}

\begin{problema}[\({\star}{\star}{\star}\)]\label{pb:14-5}
Pasá a coordenadas cartesianas geocéntricas el punto \(\varphi=-34{,}6^\circ\), \(\lambda=-58{,}4^\circ\), \(h=25\)~m. Usá \(X=(N+h)\cos\varphi\cos\lambda\), \(Y=(N+h)\cos\varphi\sin\lambda\), \(Z=(N(1-e^2)+h)\sin\varphi\).
\end{problema}
